\section{Related Works}

​​As artificial intelligence (AI) continues to expand its influence in medicine, the prediction of chronic conditions such as diabetes has become a prominent research focus. Numerous models and techniques have been proposed to address this challenge.
Ayon and Islam (2019), in Diabetes Prediction: A Deep Learning Approach, implemented deep neural networks on the PIMA Indian dataset and reported an impressive accuracy of 98.35\%. Building on this, Shams et al. (2025) introduced a novel RFE-GRU model that combines Recursive Feature Elimination with Gated Recurrent Units, achieving a notable 90.5\% accuracy. Although slightly lower than Ayon et al., their approach emphasized feature reduction and model efficiency.
Other researchers have focused on enhancing model performance through data manipulation and a range of strategies. García-Ordás et al. (2021) addressed class imbalance by applying oversampling and feature augmentation, which improved deep learning model robustness. In a broader review, Mahajan et al. (2023) explored ensemble learning techniques in other disease prediction as well, concluding that combining methods often outperform single-model approaches by leveraging the strengths of multiple classifiers.
Earlier foundational work by Jahangir et al. (2017) implemented an expert system using an auto-tuned multi-layer perceptron, demonstrating the viability of neural networks in clinical prediction tasks. Meanwhile, Khokhar et al. (2025) conducted a systematic review of AI applications in diabetes prediction, identifying key advancements, limitations, and areas requiring further investigation.
Together, these studies reflect a progression from simple neural models to complex hybrid and ensemble systems, highlighting the diverse strategies researchers have employed to tackle diabetes prediction using AI. However, determining the correct application of ensembles of different systems proves to provide further areas of research in diabetes prediction. 
